\documentclass{article}

\usepackage{amsmath, amsfonts, amssymb, amsthm}
\usepackage{tikz}
\usepackage{enumitem}
\usepackage{parskip}
\usetikzlibrary{patterns}

\theoremstyle{plain}
\newtheorem{theorem}{Theorem}[section]
\newtheorem{proposition}[theorem]{Proposition}
\newtheorem{lemma}[theorem]{Lemma}
\newtheorem{definition}[theorem]{Definition}

\title{T-tetromino tiling improvements for some classes of rectangles}
\author{Jack Grahl}

\begin{document}
\maketitle

\begin{abstract}
Hochberg showed in \cite{hochberg} that no more than 9 monominos are ever needed to tile a sufficiently large rectangle with T-tetrominos and monominos, and conjectured that no more than 5 monominos are needed for any rectangle, apart from those whose widths belong to a finite number of exceptions.
The author previously showed that this conjecture holds for all squares.
In the present work, we show tilings with the best possible number of monominos for some classes of rectangles, including all rectangles with both side lengths odd.

\end{abstract}

\section{Odd by odd rectangles}

\begin{lemma}\label{3_times_1}
    Every rectangle with dimensions $4l + 1 \times 4m + 3$ can be tiled with exactly 3 monominoes, and 3 are always needed.
    \end{lemma}
    Of course it is not possible to tile this rectangle with fewer than 3 monominoes.
    
    Suppose that $K = 4k + 1$ and $L = 4l + 3$, where $K$ and $L$ are positive integers, at least 13. (The condition mod 4 means that $L$ will be at least 15).
    
    First we find $0 < k' \leq k$ and $0 < l' \leq l$, such that we can construct a tiling of the $4k' + 1 \times 4l' + 4$ rectangle with only 3 monominos. We need $k - k'$ to be equal to $l - l'$.
    
    In other words, we are looking for $0 < K' \leq K$ and $0 < L' \leq L$ such that $K - K'$ and $L - L'$ are both the same nonnegative multiple of 4.
    
    Suppose that $K < L + 10$. This means that $K \leq L + 6$. Choose $K' = 13$ and $L' = L - K + 13 \geq 7$. Theorem 7 of \cite{hochberg} gives us a tiling of this rectangle using exactly 3 monominos.
    
    Suppose that this is not true, ie. that $K \geq L + 10$. Then take $L' = 15$. This means that $K'$ will be $K - L + 15 \geq 25$. Again, since $K' > 5$, Theorem 7 of \cite{hochberg} gives a tiling for this rectangle which uses exactly 3 monominos.
    
    Then we extend such tilings to the rectangle $4k + 1 \times 4l + 3$. Figure \ref{extending} shows how a single step works, adding 4 to both dimensions of the rectangle, without requiring any additional monominos. We apply this step $k - k' = l - l'$ times to obtain a tiling using only 3 monominos.
    
    \begin{figure}
    \input{extending_image}
    \caption{A tiling of a rectangle of dimensions $4l + 1 \times 4m + 3$ can be extended to a tiling of the rectangle $4l + 5 \times 4m + 7$, without any additional monominos.}
    \label{extending}
    \end{figure}

\begin{lemma}\label{3_times_1_strip}
    Every rectangle with dimensions $4l + 1 \times 4m + 1$ or $4l + 3 \times 4m + 3$ can be tiled with exactly 5 monominos, and 5 are always needed.
\end{lemma}

A strip of width $2$ and of length $2n+1$ can trivially be tiled with $2n$ tetrominos and $2$ monominos.
See Figure \ref{strips} for a simple visual demonstration of this.
This tiling can be used to extend a tiling of the form given in \ref{3_times_1}, adding a $2$ to one or other dimension, and $2$ monominoes.

\begin{figure}
\input{strips_image}
\caption{Tilings of strips of width 2 and length $4n + 1$ or and $4n+3$, with 2 monominos.}
\label{strips}
\end{figure}

That rectangles of this form cannot be tiled with less than $5$ monominoes, is the main result of \cite{zhan}.


\section{Remaining cases to be resolved}
Hochberg divided all rectangles into classes based on their side lengths modulo $4$, and showed the best known construction for each class.
He identified five classes where the best known construction was not known to be optimal:
\begin{itemize}
	\item $4k + 1 \times 4m + 3$: Hochberg showed a construction with $7$ monominoes, concluding that the optimal number is either $3$ or $7$. Lemma \ref{3_times_1} shows that $3$ is always possible, resolving this case.  
	\item $4k + 1 \times 4m + 1$: Zhan proved that $1$ is not possible. Hochberg showed a construction with $9$, concluding that the optimal number is either $5$ or $9$. Lemma \ref{3_times_1_strip} shows that $5$ is always possible, resolving this case.  
	\item $4k + 3 \times 4m + 3$: Zhan proved that $1$ is not possible. Hochberg showed a construction with $9$, concluding that the optimal number is either $5$ or $9$. Lemma \ref{3_times_1_strip} shows that $5$ is always possible, resolving this case.  
	\item $4k + 2 \times 4m + 1$: Hochberg showed a construction with $6$, concluding that either 2 or 6 is optimal. We show constructions with only $2$ monominoes for some special cases in the next section.
	\item $4k + 2 \times 4m + 3$: Hochberg showed a construction with $6$, concluding that either 2 or 6 is optimal. We show constructions with only $2$ monominoes for some special cases in the next section.
\end{itemize}

The resolution of the first 3 cases here tells us that, in the terminology used by Hochberg, the \emph{gap cap} of the T-tetromino is either 5 or 6.
To finally confirm its exact value, we need to establish one of the following complementary possibilities, which treat the latter two cases above:

\begin{enumerate}[label=\textnormal{\Roman*}]
	\item \label{gap_cap_5} For some $N \in \mathbb{N}$, there is a way of tiling all rectangles with side lengths $4p + 2$ and $2q + 1$, with both side lengths greater than $N$, using only $2$ monominoes.
	\item \label{gap_cap_6} There is a family of rectangles with side lengths $4p + 2$ and $2q + 1$, containing rectangles with both sides arbitrarily large, such that none can be tiled using less than $6$ monominoes.
\end{enumerate}

Note that a proof of \ref{gap_cap_6} can not just establish that $6$ monominoes are necessary for any infinite family of rectangles.
It must contain rectangles, with arbitrarily large numbers for both dimensions.
The possibility of a family of counterexamples where one of two dimensions is bounded is explicitly excluded by the definition of the gap cap.

Hochberg's conjecture that the gap cap is $5$ is equivalent (given the other results above) to \ref{gap_cap_5}.
In the next section, we give constructions using only $2$ monominoes for certain families of rectangles.
This provides `evidence' for the conjecture, however there is no clear way to extend these constructions to other rectangle sizes.

For clarity we summarize the best known results by modulo classes in Table \ref{rect_tab}.

\begin{table}
    \input{rectangle_table}
    \caption{Minimal required numbers of monominoes.}\label{rect_tab}
    \end{table}

\section{Odd by even rectangles}

We would like to show a construction which works for any rectangle $m times n$, where $m$ is an odd number and $n$ is a number of the form $4k + 2$, using t-tetrominoes and 2 monominoes.
We don't have such a construction, but we have a construction like this for a couple of families of rectangles.

\subsection{$2n \times n$ rectangles}

\begin{lemma}\label{2_times_1}
Suppose that $n = 2k+1$ is an odd number. Then a rectangle $n \times 2n$ can be tiled with 2 monominoes.  
\end{lemma}

We use the fact, from \cite{grahl_sq}, that for any odd number $n$, we can tile the 'nibbled square', the $n \times n$ square with 3 cells removed at one corner, and one cell removed from an adjacent corner, with a single monomino.
Figure \ref{nibbled} shows what this shape looks like.

\begin{figure}
\input{nibbled_square}
\caption{The `nibbled square of size $n$', which is missing the lattice squares at $(0, 0), (1, 0), (0, 1)$ and at $(0, n-1)$}
\label{nibbled}
\end{figure}

Figure \ref{n_times_2n} shows that we can put together two copies of the nibbled square of size $n$, one reflected, plus two additional t-tetrominos, to tile the $2n \times n$ rectangle.
Since each nibbled square is tiled with t-tetrominoes and a single monomino, the $2n \times n$ rectangle is tiled with only $2$ monominoes.

\begin{figure}
\input{n_times_2n_image}
\caption{Tiling a rectangle $n \times 2n$, where $n$ is an odd number.}
\label{n_times_2n}
\end{figure}


\subsection{$5n \times 2n$ rectangles}

\begin{lemma}\label{5_times_2}  
Suppose that $n = 2k+1$ is an odd number. Then a rectangle $5n \times 2n$ can be tiled with 2 monominoes.  
\end{lemma}

\begin{figure}
\input{2n_times_5n_image}
\caption{Tiling a rectangle $2n \times 5n$, where $n$ is an odd number.}
\label{2n_times_5n}
\end{figure}

\bibliography{polyomino}{}
\bibliographystyle{plain}
\end{document}
